% Packing k^2 - ceil(sqrt(32k)) - 1 unit squares in a square of side less than k
% Kevin Fang, 2026-10-24. LaTeX source of the note built from docs/theorem-a.md (v0.3).
% Build: tectonic main.tex   (or: pdflatex main.tex; pdflatex main.tex)
\documentclass[11pt,letterpaper]{amsart}

\usepackage[T1]{fontenc}
\usepackage{lmodern}
\usepackage{amsmath,amssymb}
\usepackage{booktabs}
\usepackage[margin=1in]{geometry}
\usepackage{xurl}
\usepackage{xcolor}
\usepackage[colorlinks=true,linkcolor=blue!60!black,citecolor=blue!60!black,urlcolor=blue!60!black]{hyperref}

\DeclareUrlCommand\hash{\urlstyle{rm}}
\hypersetup{pdftitle={Packing k\textasciicircum 2 - ceil(sqrt(32k)) - 1 unit squares in a square of side less than k},pdfauthor={Kevin Fang}}
\emergencystretch=1.5em

\theoremstyle{plain}
\newtheorem{theorem}{Theorem}
\newtheorem{lemma}[theorem]{Lemma}
\theoremstyle{remark}
\newtheorem*{remark}{Remark}

\newcommand{\cs}{c^{*}}
\newcommand{\fl}[1]{\lfloor #1 \rfloor}

\begin{document}

\title[Packing $k^2-\lceil\sqrt{32k}\,\rceil-1$ unit squares below side $k$]{Packing $k^2-\lceil\sqrt{32k}\,\rceil-1$ unit squares in a square of side less than $k$}
\author{Kevin Fang}
\email{me@kevinfang.tech}
\date{2026-10-24}
\thanks{Draft v0.2, 2026-10-24, doi:10.5281/zenodo.23624565; v0.2 credits the earlier result of Daniel \cite{Dan26b} (v0.1: doi:10.5281/zenodo.23624412). Code and certificates: \cite{Fan26b}, folder \texttt{theorem-a/}, commit \hash{b1b3f6b}.}

\begin{abstract}
Let $s(N)$ be the side of the smallest square containing $N$ unit squares with disjoint interiors, and let $\cs(k)$ be the largest $c$ with $s(k^2-c)=k$. We prove $\cs(k)<4\sqrt{2k}+1$ for every integer $k\ge6$, that is, $s(k^2-\lceil\sqrt{32k}\,\rceil-1)<k$. The construction is Ryu's L-shaped packing: an aligned grid and two bands of $b$ tilted rows. Both ends of each band are filled with axis-parallel rows, and $b$ is the odd integer nearest $\sqrt{2k}$. The odd width keeps the rounding loss of each filled end below $3/2$. The same idea and the same constant $4\sqrt2$ were found earlier by Evan Daniel, who proves $\cs(k)\le4\lceil\sqrt{2k-7}\,\rceil-1$ for $k\ge7$ in Lean; for $k\ge89$ the two integer bounds differ by at most $2$. This note is an independent second proof, with a different split of the free length and an odd band width. Exact rational certificates for $k\le600$ are checked by two independently written programs.
\end{abstract}

\maketitle

\section{Introduction}

For $x>0$ let $N(x)$ be the largest number of unit squares that fit, with pairwise disjoint interiors and rotations allowed, in a square of side $x$. For an integer $k$, $N(k)=k^2$. Let
\[
M(k)=\max\{N : s(N)<k\},\qquad w(k)=k^2-M(k),\qquad \cs(k)=\max\{c : s(k^2-c)=k\}.
\]
Then $s(k^2-c)<k$ exactly when $M(k)\ge k^2-c$, so $\cs(k)=w(k)-1$.

Erd\H{o}s and Graham \cite{EG75} showed that the waste $x^2-N(x)$ is $O(x^{7/11})$, which makes $w(k)$ sublinear. The best exponent is now $3/5$ \cite{Bui25,McC26} (the earlier claim \cite{CG20} contained an error \cite{AB26}). For integer sides approached from below, Ryu \cite{Ryu26} proved, among other things:
\begin{itemize}
\item $\cs(k)\le 8\lceil\sqrt{k-4}\,\rceil-1$ for every $k\ge6$ (his Theorem~1.4), from an explicit L-shaped packing (his Theorem~1.3);
\item $\cs(k)<20.668\,k^{2/5}+0.623$ for $k\ge1.6\cdot10^{7}$ (Theorem~1.2), and bounds of order $k^{3/8}$ from $k\ge2.48\cdot10^{12}$ (Theorem~1.1);
\item computer-assisted bounds at five values of $k$ between $10^5$ and $10^8$ (Theorem~1.6).
\end{itemize}
Daniel \cite{Dan26b} proved in Lean, on 18 October 2026, $\cs(k)\le4\lceil\sqrt{2k-7}\,\rceil-1$ for every $k\ge7$, by filling the band ends of Ryu's L with axis-parallel layers, and $\cs(k)<18.9\,k^{2/5}+35.7\,k^{1/5}+122$ for $k\ge3\cdot10^5$ from Ryu's staircase ends. We learned of this after writing v0.1 of this note; Theorem~\ref{thm:A} below is an independent proof of a bound of the same strength. Lower bounds of order $\log k$, $k^{1/4}$ and $k^{1/3}$ for large $k$ are also due to Ryu \cite{Ryu26b}. Arslanov, Mustafin and Shangitbayev \cite{AMS21} proved $\cs(k)\le k-1$ for $k\ge12$. In an earlier note \cite{Fan26a} we proved $\cs(k)\le k-1-F(k)$ for $k\ge15$, with $F(k)=1+\max\{p+2q:16p+21q\le k-15\}$.

\begin{theorem}\label{thm:A}
For every integer $k\ge6$, $w(k)<4\sqrt{2k}+2$. Equivalently $\cs(k)<4\sqrt{2k}+1$, so $\cs(k)\le\lceil\sqrt{32k}\,\rceil$, and $s(n)<k$ for every integer $n\le k^2-\lceil\sqrt{32k}\,\rceil-1$.
\end{theorem}

The packing is Ryu's L (his Sections~3--4), and we use his end-region framework (his Section~6): lifted band ends, the far end as the $180^\circ$ image of the near one, and end fills assembled with the rows. Daniel \cite{Dan26b} fills the band ends with axis-parallel layers in the same way. What differs here is the split of the free length (lift $0$ at one end of each band, all of it at the other), the odd band width, and the estimate of Lemma~\ref{lem:fill}, which together give the bound in the form $w(k)<4\sqrt{2k}+2$. To keep the note self-contained, we prove everything we use.

\section{The band}\label{sec:band}

Fix an integer $b\ge3$ and a rational $t=\tan(\theta/2)>1/b$. Write $c=\cos\theta=(1-t^2)/(1+t^2)$, $s=\sin\theta=2t/(1+t^2)$, $T=s/c$, $\tau=c/s$, $\sec=1/c$, and
\[
W=bc+s\quad(\text{band width}),\qquad h=bs+c\quad(\text{row height}).
\]
At $t_0=1/b$: $c_0=(b^2-1)/(b^2+1)$, $s_0=2b/(b^2+1)$, $\tau_0=(b^2-1)/(2b)$, $W_0=b$, $h_0=(3b^2-1)/(b^2+1)<3$. Since $dW/d\theta=c-bs<0$ for $\tan\theta>1/b$, and $\theta_0=2\arctan(1/b)>\arctan(1/b)$, we have $W<b$ for every $t>1/b$.

Use band coordinates $X\in[0,W]$ (across) and $Y\in[0,L]$ (along). With $(u,v)$ the coordinates along $(c,s)$ and $(-s,c)$, the rectangle $P=\{0\le u\le b,\ 0\le v\le1\}$ is the union of the $b$ unit squares $\{i\le u\le i+1\}$. \emph{Row $j$} is $P$ translated so that its corner $u=v=0$ is at $(s,\,y_0+j\sec)$, for $0\le j<R$.

\emph{(R1)} A point of $P$ has $X-s=uc-vs\in[-s,bc]$, so every row lies in $0\le X\le W$. \emph{(R2)} The shift $(0,\sec)$ changes $v$ by exactly $1$, so in row~$0$'s frame row $j$ lies in $\{j\le v\le j+1\}$ and the rows have disjoint interiors. \emph{(R3)} The lowest point of row $j$ is its corner, at height $y_0+j\sec$, and its highest is at height $y_0+j\sec+h$. So the rows fit in $[0,L]$ when $y_0+(R-1)\sec+h\le L$. Let $R=R(L)=\fl{(L-h)/\sec}+1$, the most rows with $y_0=0$, and let $F=L-(R-1)\sec-h\in[0,\sec)$ be the free length.

\emph{End regions.} With row~$0$ at lift $y$, let $\ell(X)=y+(X-s)T$ for $X\ge s$ and $\ell(X)=y+(s-X)/T$ for $X\le s$. \emph{(E1)} Every row lies in $\{Y\ge\ell(X)\}$: in row~$0$'s frame a point of any row has $u,v\ge0$, and $Y-y-(X-s)T=vc+vs^2/c\ge0$ for $X\ge s$, while $Y-y-(s-X)/T=us+uc^2/s\ge0$ for $X\le s$. So a unit square in $\{0\le X\le W,\ 0\le Y\le\ell(X)\}$ meets no row in an interior point.

\emph{Fill.} Let $a_i(y)=\max(0,\,s+(i+1-y)\tau)$ and $m_i(y)=\max(0,\fl{W-a_i(y)})$. Level $i$ holds the $m_i$ unit squares $[W-m_i+q,\,W-m_i+q+1]\times[i,i+1]$, $0\le q<m_i$, and $\mathrm{fill}(y)=\sum_i m_i(y)$. \emph{(E2)} These squares lie under $\ell$. If $a_i\ge s$, $\ell$ increases on $[a_i,W]$ and $\ell(a_i)\ge i+1$. If $a_i<s$, then $i+1<y\le\min\ell$.

\emph{(E3)} The $180^\circ$ rotation about $(W/2,L/2)$ maps the rows onto rows with lifts $y_1+(R-1-j)\sec$, where $y_1=L-y_0-(R-1)\sec-h$. So the region over the last row is a rotated end region with lift $y_1$, and takes the rotated fill. The two fills are disjoint, since each vertical line $X=\mathrm{const}$ meets row~$0$ and row $R-1$. A band of length $L$ therefore holds
\[
C(L)=bR+\mathrm{fill}(y_0)+\mathrm{fill}(y_1)
\]
unit squares. Below, $y_0=0$ and $y_1=F$.

\section{The packing}

Let $S=k-b+W<k$, and suppose $k-b>h$. Put three pieces in $[0,S]^2$:
\begin{itemize}
\item the aligned grid $[0,k-b]^2$;
\item the band of length $S$ in $[k-b,S]\times[0,S]$ (with $x=k-b+X$, $y=Y$);
\item the band of length $k-b$, reflected in $y=x$, in $[0,k-b]\times[k-b,S]$.
\end{itemize}
They have disjoint interiors. Since $k^2-(k-b)^2=bk+b(k-b)$,
\begin{equation}\label{eq:waste}
w(k)\le b(k-S)+\bigl[bS-C(S)\bigr]+\bigl[b(k-b)-C(k-b)\bigr].
\end{equation}

\section{The end lemma}

From now on $b$ is odd. Put $\eta=b/2-\tau$; at $t_0$, $\eta_0=1/(2b)$. Say $t>1/b$ is \emph{close} if
\begin{itemize}
\item[(a)] $\fl{W-s-\tau}=(b-1)/2$;
\item[(b)] $0\le\eta<1/2$;
\item[(c)] $2\tau-bc\ge0$;
\item[(d)] $h<3$.
\end{itemize}
At $t_0$ we have $W_0-s_0-\tau_0=(b-1)/2+\delta_b$ with $\delta_b=\tfrac12+\tfrac1{2b}-\tfrac{2b}{b^2+1}\in(0,1)$; here $\delta_b>0$ because $b^3-3b^2+b+1>0$ for $b\ge3$. Also $2\tau_0-bc_0=(b^2-1)/(b(b^2+1))>0$. So (a)--(d) hold strictly at $t_0$, and by continuity on an interval $(1/b,1/b+\varepsilon_b)$.

\begin{lemma}\label{lem:fill}
If $t$ is close, then $\mathrm{fill}(0)\ge(b-1)/2$, and for $0\le y<\sec$,
\[
\mathrm{fill}(y)>\tfrac{b-1}{2}+2z-1-2\eta,\qquad z=W-s-2\tau+y\tau.
\]
\end{lemma}

\begin{proof}
$\mathrm{fill}(0)\ge m_0(0)=\fl{W-s-\tau}=(b-1)/2$ by (a). For $0\le y<\sec$ and $i\ge0$, the clipping in $a_i$ is inactive:
\[
s+(i+1-y)\tau>s+(1-\sec)\tau=c(1-c)/s\ge0.
\]
So level~$0$ holds at least $\fl{z+\tau}$ squares and level~$1$ at least $\fl z$. Since $\tau=\frac{b-1}2+\frac12-\eta$ with $\frac{b-1}2\in\mathbb Z$, together they hold at least $\frac{b-1}2+\fl{z+\frac12-\eta}+\fl z$. For $0\le\eta<\frac12$ and real $z$ with fractional part $f$, the sum $\fl z+\fl{z+\frac12-\eta}$ is $2z-2f>2z-1-2\eta$ if $f<\frac12+\eta$, and $2z-2f+1>2z-1$ otherwise.
\end{proof}

The odd width matters: it puts the arguments of the two floors half a unit apart. For even $b$ the arguments differ by an integer minus $1/(2b)$, and this pairing gives nothing.

\begin{lemma}\label{lem:band}
If $t$ is close and $L>h$, then with $R=R(L)$, $y_0=0$ and $y_1=F$,
\begin{equation}\label{eq:band}
bL-C(L)<bL(1-c)+bch-2b+2-2(W-s-2\tau)+2\eta.
\end{equation}
At $t_0$ the right side equals $2bL/(b^2+1)+b+2-\kappa_b$, where $\kappa_b=(7b^4+1)/\bigl(b(b^2+1)^2\bigr)$.
\end{lemma}

\begin{proof}
$bR=b+bc(L-h-F)$. By Lemma~\ref{lem:fill}, $C(L)>bR+(b-1)+2(W-s-2\tau)+2F\tau-1-2\eta$. The terms in $F$ sum to $F(2\tau-bc)\ge0$ by (c), so drop them and rearrange. The value at $t_0$ uses $1-c_0=2/(b^2+1)$, $bc_0h_0=3b-b(10b^2+2)/(b^2+1)^2$, $W_0-s_0-2\tau_0=-(b^2-1)/(b(b^2+1))$ and $2\eta_0=1/b$.
\end{proof}

\section{Proof of Theorem~\ref{thm:A}}

Let $k\ge6$ and let $b$ be an odd integer with $|b-\sqrt{2k}|\le1$. Then $b\ge3$. For $k\ge8$, $k-b\ge k-\sqrt{2k}-1\ge3>h$ by (d). For $k=6,7$ the only choice is $b=3$, and $k-b\ge3>h$. For close $t$, \eqref{eq:waste} and \eqref{eq:band} with $L=S$ and $L=k-b$ give $k^2-N(t)<H(t)$, where $N(t)$ is the number of squares packed and $H(t)$ is $b(k-S)$ plus the two right sides of \eqref{eq:band}.

This holds for each $t$ separately, even though $R$, $F$ and the floors jump with $t$. $H$ is continuous, and as $t\to1/b^+$, $S\to k$ and
\[
H(t)\to B(k)=\frac{4kb}{b^2+1}+2b+2+\frac{2}{b^2+1}-2\kappa_b.
\]
Now $4kb/(b^2+1)<4k/b$, and $4k/b+2b=4\sqrt{2k}+\frac2b(b-\sqrt{2k})^2\le4\sqrt{2k}+\frac2b$. So
\[
B(k)<4\sqrt{2k}+2-2\Bigl(\kappa_b-\frac1b-\frac1{b^2+1}\Bigr),
\]
and $b(b^2+1)^2\bigl(\kappa_b-\frac1b-\frac1{b^2+1}\bigr)=6b^4-b^3-2b^2-b>0$. Hence $B(k)<4\sqrt{2k}+2$. Some close $t$ has $H(t)<4\sqrt{2k}+2$, and it gives a packing with $S<k$ and $k^2-N(t)<4\sqrt{2k}+2$. \qed

\begin{remark}[Sharpness]
For odd $b\ge5$ and $k=(b^2-b+8)/2$, the construction at close $t$ has waste exactly $4b$, which is $4\sqrt{2k}+2-31/(2b)+O(b^{-2})$; for example $k=59$, $b=11$, $w=44$. So the additive constant $2$ cannot be lowered for this construction.
\end{remark}

\section{Checks and certificates}

The construction is implemented in exact rational arithmetic (\texttt{lband.py} in \cite{Fan26b}).
\begin{itemize}
\item \emph{Admissibility.} For each $k$ the program chooses the squeeze $k-S$ as $1/(10^6k^3)$, divided by $10$ until (a)--(d) and $H(t)<4\sqrt{2k}+2$ hold. It decides those conditions exactly, so every layout it returns satisfies the theorem's bound by Lemma~\ref{lem:band}.
\item \emph{A fixed squeeze is not enough.} With a fixed squeeze of $10^{-7}$ the same recipe has waste $1178=4\sqrt{2k}+2$ at $k=43218$, which violates the strict bound.
\item \emph{Every step checked.} Every inequality in the proof was checked exactly for $6\le k\le20000$, and the bound for $6\le k\le60000$. The least slack $B(k)-w$ is $8019401/78023840599$, at $k=19705$.
\item \emph{Certificates.} Strict certificates (every wall and pair margin positive) of the construction for $6\le k\le40$ and $k=101,300,600$ pass our checker and Evan Daniel's \cite{Dan26}. The largest has $359860$ squares. A program that is not all-pairs runs both checkers unmodified on overlapping $4\times4$ windows: two unit squares whose centres are more than $\sqrt2$ apart cannot overlap, and both checkers accept such pairs without a separating-axis test.
\end{itemize}

The construction is not optimized. Searching over $b$, the number of rows and the split of the free length between the two ends gives waste at most $4\sqrt{2k}+1.18$ for $15\le k\le300$; for example $9946$ squares below side $100$. A bounded search for better end and corner pieces by local optimization found gains of one or two squares only (at $k=25,30,40$), not growing with $b$. That is consistent with the heuristic that a boundary between the aligned and the tilted lattice wastes about $1/2$ per unit length.

\section{Comparison}

The integer form of the theorem is $\cs(k)\le\lceil\sqrt{32k}\,\rceil$.
\begin{itemize}
\item \emph{Daniel's bound} \cite{Dan26b}, $4\lceil\sqrt{2k-7}\,\rceil-1$. For $7\le k<10^6$ ours is smaller at $495210$ values of $k$, his at $255215$, and they tie at the remaining $249568$. For $k\ge89$ they never differ by more than $2$; at $17$ values of $k$ below $89$ his is smaller by $3$ to $5$. His $k^{2/5}$ bound is smaller from about $k=7.4\cdot10^5$.
\item \emph{Ryu's Theorem~1.4.} Ours is smaller at $k=6$ and for every $k\ge9$; they tie at $k=7$, and Ryu's is smaller at $k=8$.
\item \emph{Ryu's Theorem~1.3}, optimized over $b$. Ours is smaller for $k\ge14$.
\item \emph{\cite{Fan26a}.} Ours is smaller from $k=42$ (tie at $40$, $41$), and larger for $15\le k\le39$.
\item \emph{Superseded} by Ryu's Theorem~1.2 for $k\ge1.6\cdot10^7$, by his Theorem~1.1 tiers beyond that, and by his computer-assisted values: at $k=10^5,10^6,10^7$ they give $\cs\le1583,3972,10039$, against our $1789,5657,17889$.
\end{itemize}
So Theorem~\ref{thm:A} is of the same strength as Daniel's earlier bound, not a new record for any range of $k$.

\section*{Acknowledgements and credit}

The L-shaped packing, its end-region framework and most of the bounds compared above are Sungjoon Ryu's. The axis-layer filling of the band ends and the bound $4\lceil\sqrt{2k-7}\,\rceil-1$ were found first by Evan Daniel \cite{Dan26b}, whose certificate format and checker we also use. The Squares Project (Joshua Levy) keeps the register in which these bounds are compared. The construction code, checkers and certificates were written and run with Claude (Anthropic) as a coding and research agent, directed and reviewed by the author. Two blind language-model reviews of the proof (Claude Opus; OpenAI Codex, GPT-6 Astra) found errors that were corrected. That is not an external mathematical review.

\begin{thebibliography}{AMS21}

\bibitem[AB26]{AB26} M. Z. Arslanov, H. D. Bui, Note on ``Efficient packings of unit squares in a large square'', Discrete Comput. Geom. 76 (2026), 805--812, doi:10.1007/s00454-025-00767-w.

\bibitem[AMS21]{AMS21} M. Z. Arslanov, S. A. Mustafin, Z. K. Shangitbayev, Improved packings of $n(n-1)$ unit squares in a square, Electron. J. Combin. 28(4) (2021), P4.22, doi:10.37236/8586.

\bibitem[Bui25]{Bui25} H. D. Bui, Square packing with $O(x^{0.6})$ wasted area, arXiv:2508.04603 (preprint; v2, 15 March 2026).

\bibitem[CG20]{CG20} F. Chung, R. Graham, Efficient packings of unit squares in a large square, Discrete Comput. Geom. 64 (2020), 690--699, doi:10.1007/s00454-019-00088-9.

\bibitem[Dan26]{Dan26} E. Daniel, \url{https://github.com/evand/square-packing}, \texttt{search/exact/verify\_cert.py} at commit \texttt{29d0cee87c1d717f7b4497b8fb0ee17ac9d0087d}.

\bibitem[Dan26b]{Dan26b} E. Daniel, Upper bounds on $\cs(k)$, in Lean, \url{https://github.com/evand/square-packing}, \texttt{certificates/k2-minus-c/} at commit \texttt{aa12b98} (18 October 2026); write-up \url{https://evand.github.io/square-packing/k2c/}.

\bibitem[EG75]{EG75} P. Erd\H{o}s, R. L. Graham, On packing squares with equal squares, J. Combin. Theory Ser. A 19 (1975), 119--123, doi:10.1016/0097-3165(75)90099-0.

\bibitem[Fan26a]{Fan26a} K. Fang, Packing $n^2-n+1$ unit squares in a square of side less than $n$, for every $n\ge15$, draft v0.9 (2026), doi:10.5281/zenodo.23463481.

\bibitem[Fan26b]{Fan26b} K. Fang, squarepack-certs, \url{https://github.com/TheSnakeFang/squarepack-certs}, folder \texttt{theorem-a/}, commit \hash{b1b3f6bda13739c9eeecc1ceaefec9178294e0ca}.

\bibitem[McC26]{McC26} R. McClenagan, Optimally packing a large square by unit squares, arXiv:2602.01484 (preprint, 2026).

\bibitem[Ryu26]{Ryu26} S. Ryu, Packing $k^2-c$ unit squares: an upper bound of order $k^{3/8}$ for the deficiency, preprint, version 1.1 (10 October 2026), doi:10.5281/zenodo.23278927 (version 1.0: doi:10.5281/zenodo.23256655).

\bibitem[Ryu26b]{Ryu26b} S. Ryu, Packing $k^2-c$ unit squares: $s(k^2-c)=k$ for all large $k$, preprint, version 1.2 (2026), doi:10.5281/zenodo.23194104; and the companion preprints on lower bounds of order $k^{1/4}$ (doi:10.5281/zenodo.23212643) and $k^{1/3}$ (doi:10.5281/zenodo.23213564).

\end{thebibliography}

\end{document}
